\begin{definition} \label{definition:invariant_subgroup_of_a_group_acting_on_an_abelian_group}

Let $A$ be an \CrefAndHyperrefIfExist{definition:group}{abelian group} with a \CrefAndHyperrefIfExist{definition:action_of_a_group_on_an_abelian_group}{$G$-action}, say on the left. 
The \hldef{group of $G$-invariant elements}, denoted \hl{$A^G$} (this notation is also standard when $G$ \CrefAndHyperrefIfExist{definition:action_of_a_group_scheme_on_a_scheme}{acts on} $A$ on the right), is \CrefAndHyperrefIfExist{definition:subgroup_of_a_group}{subgroup} of $A$ defined as:
$$ A^G = \{ a \in A \mid g \cdot a = a \text{ for all } g \in G \}.$$
\TextIfExists{definition:fixed_point_set_of_a_group_acting_on_a_set}{Note that this coincides with the fixed point sete of $G$ in $A$ as per \Cref{definition:fixed_point_set_of_a_group_acting_on_a_set}.}

If $A$ is an \CrefAndHyperrefIfExist{definition:left_right_module_over_a_group_ring_of_a_group_over_a_ring}{(left/right) $R[G]$-module} where $R[G]$ is the \CrefAndHyperrefIfExist{definition:group_ring_of_a_group_over_a_ring}{group ring}, then $A^G$ is an (left/right) \CrefAndHyperrefIfExist{definition:module_of_a_ring}{$R$-module}. In fact, $A^G$ is naturally regarded as a \CrefAndHyperrefIfExist{definition:trivial_G_module_of_a_group_over_a_ring}{trivial $R[G]$-module}.
\end{definition}